\documentclass[11pt]{article}
\usepackage{amsmath,amssymb,amsthm}
\usepackage[margin=1.1in]{geometry}
\newtheorem{theorem}{Theorem}
\newtheorem{lemma}[theorem]{Lemma}
\newtheorem{corollary}[theorem]{Corollary}
\newtheorem{proposition}[theorem]{Proposition}
\newtheorem{remark}[theorem]{Remark}
\newcommand{\Sym}{\operatorname{Sym}}
\newcommand{\Gm}{\Gamma_{\mathbb C}}
\title{Transport closure:\\ direct sums of weight systems compose exactly}
\author{Samuel Lavery\thanks{Machine-checked companion:
\texttt{RequestProject/TransportClosure.lean}, 19 named declarations, each at axiom footprint
$\{\texttt{propext},\texttt{Classical.choice},\texttt{Quot.sound}\}$ or a subset, no
\texttt{sorry}, no \texttt{axiom}, zero \texttt{sorryAx}.  Every argument below is given in the
text; the Lean names locate the corresponding formal statement.}}
\date{Draft --- \today}

\begin{document}
\maketitle

\begin{abstract}
A single functorial lift is one theorem; a transport calculus is a system.  This note proves
the closure property of the carrier-transport machinery: when two weight systems are combined
by direct sum, every layer of the arithmetic passport composes exactly --- local Euler
coefficients by Cauchy product, global coefficient banks by Dirichlet convolution, prescribed
completions by multiplication at the product conductor and concatenated shift list, and
self-dual functional equations by multiplication.  The abstract composition laws (identity,
associativity, faithfulness under composition, base-change degrees multiplying) are compiled
separately.  Two arithmetic instances tie the calculus to a live eigenform: the tensor-square
decomposition $\mathrm{std}\otimes\mathrm{std}=\Sym^2\oplus\mathbf 1$, whose bank is the rank-4
Rankin bank $\zeta\star\mathrm{sym2Bank}$, and the plethysm
$\Sym^2\!\circ\Sym^2=\Sym^4\oplus\mathbf 1$, whose bank is $\zeta\star(\Sym^4\text{-bank})$ ---
both derived from the closure law, not verified numerically.  Everything is unconditional and
machine-checked.
\end{abstract}

\section{Setup}

A \emph{weight system} of rank $d$ is a family $w=(w_i)_{i\in I}$ of $d$ nonzero complex
numbers, one family per prime in the global setting.  Its \emph{local Euler coefficient} at
level $j$ is the complete homogeneous symmetric function
\[
  h_j(w)\;=\;\sum_{\substack{(l_i)\ge0\\ \sum l_i=j}}\ \prod_{i\in I}w_i^{\,l_i},
\]
the $p^j$-coefficient of $\prod_i(1-w_ip^{-s})^{-1}$.  Its \emph{global bank} is the
multiplicative arithmetic function assembled from the local coefficients over the prime
factorization,
$c_w(n)=\prod_{p\mid n}h_{v_p(n)}(w_p)$; formally the shift-by-one of the all-place
coefficient sequence (\texttt{coefficientArithmetic} of
\texttt{radialGlobalSatakeCoeff}).  Its \emph{prescribed completion} at conductor $C>0$ and
shift list $(\mu_j)$ is $C^s\prod_j\Gm(s+\mu_j)\cdot\sum_n c_w(n)n^{-s}$.

Two conventions used throughout: banks are compared as arithmetic functions (so ``the bank of
$w$'' means the shifted sequence), and a \emph{direct sum} $u\oplus v$ of weight systems is the
disjoint union of the two families at every prime.

\section{The closure laws}

\begin{theorem}[local closure: the Cauchy law]\label{thm:cauchy}
For weight systems $u$ of rank $d_1$ and $v$ of rank $d_2$ and every level $j$,
\[
  h_j(u\oplus v)\;=\;\sum_{a+b=j}h_a(u)\,h_b(v).
\]
\end{theorem}
\begin{proof}
Fiber the exponent vectors of the union over the total exponent $a$ carried by the
$u$-channels.  The slice with $u$-total $a$ is the product of the $u$-antidiagonal at level
$a$ and the $v$-antidiagonal at level $j-a$, and the monomial factors accordingly:
$\prod_{i}(u\oplus v)_i^{l_i}=\bigl(\prod u_i^{l_i}\bigr)\bigl(\prod v_i^{l_i}\bigr)$.
Summing slices gives the Cauchy product.  The bijection is
$l\mapsto(a,\,l|_u,\,l|_v)$ with inverse gluing, and both roundtrips are identities
coordinate by coordinate.  (\texttt{TransportClosure.radialLocalEulerCoeff\_sumElim}.)
\end{proof}

Two degenerate values calibrate the law: the local coefficient is invariant under relabeling
of channels (\texttt{radialLocalEulerCoeff\_equiv}), and the unit system --- one channel of
weight $1$ --- has $h_j=1$ for every $j$ (\texttt{radialLocalEulerCoeff\_unit}), i.e.\ its bank
is $\zeta$.

\begin{theorem}[global closure: banks convolve]\label{thm:bank}
The global bank of any weight system is multiplicative
(\textup{\texttt{isMultiplicative\_bankArithmetic}}), its prime-power values are the local
coefficients (\textup{\texttt{bankArithmetic\_prime\_pow}}), and the bank of a direct sum is
the Dirichlet convolution of the banks:
\[
  c_{u\oplus v}\;=\;c_u\star c_v .
\]
\end{theorem}
\begin{proof}
Multiplicativity: at coprime $m,n$ the prime factors split and the factorization exponents
restrict, so the defining product splits.  Both sides of the convolution identity are then
multiplicative, so they agree once they agree at prime powers; there the convolution is the
Cauchy product of prime-power values, $(c_u\star c_v)(p^j)=\sum_{a+b=j}c_u(p^a)c_v(p^b)$
(divisors of $p^j$ are the $p^a$; \texttt{mul\_apply\_prime\_pow}), which is
Theorem~\ref{thm:cauchy}.  (\texttt{bankArithmetic\_sumElim}.)
\end{proof}

\begin{theorem}[completion and functional-equation closure]\label{thm:completion}
Polynomial Satake dual pairs are closed under direct sum
(\textup{\texttt{sumPair}}; exponent the maximum of the two), with coefficient closure
\textup{(\texttt{sumPair\_primalCoeff})}.  On any common point of absolute convergence the
prescribed completion of the sum, at the product conductor and the concatenated shift list, is
the product of the two prescribed completions:
\[
  (C_1C_2)^s\!\!\prod_{\mu\in\mu_1\mathbin{+\!\!+}\mu_2}\!\!\Gm(s+\mu)\cdot L(c_{u\oplus v},s)
  \;=\;\Bigl(C_1^s\prod_{\mu_1}\Gm\cdot L(c_u,s)\Bigr)
        \Bigl(C_2^s\prod_{\mu_2}\Gm\cdot L(c_v,s)\Bigr),
\]
and self-dual functional equations multiply: if $F(1-s)=F(s)$ and $G(1-s)=G(s)$ then
$(FG)(1-s)=(FG)(s)$.
\end{theorem}
\begin{proof}
The conductor factor splits since both conductors are nonnegative reals; the $\Gamma$-product
splits over the concatenation; and $L(c_u\star c_v,s)=L(c_u,s)L(c_v,s)$ where both factors
converge absolutely (Dirichlet-series multiplication, with the shift-by-one bookkeeping
handled by the bridge \texttt{dirichlet\_shift\_eq\_LSeries}).  The functional-equation clause
is immediate.  (\texttt{completedReadout\_sumPair}, \texttt{FE\_mul}.)
\end{proof}

\begin{remark}[the abstract layer]
Identity, associativity, closure of faithfulness under composition, and the base-change
instance $\mathrm{bc}_p\circ\mathrm{bc}_q=\mathrm{bc}_{pq}$ are compiled separately
(\texttt{CarrierFunctoriality.faithful\_comp}, \texttt{comp\_assoc}, \texttt{id\_comp},
\texttt{comp\_id}, \texttt{bc\_comp}, \texttt{bc\_one}).  Together with
Theorems~\ref{thm:cauchy}--\ref{thm:completion} the transports form a composition calculus in
which every layer of the arithmetic passport is preserved --- the systems property, as opposed
to any single lift.
\end{remark}

\section{Two live instances at an eigenform}

Let $f$ be a level-one holomorphic Hecke eigenform of weight $k$ with Satake parameters
$\alpha_p$, and recall from the $\Sym^2$ companion the clock systems
$\Sym^2\!:(\alpha^2,1,\alpha^{-2})$ and $\Sym^4\!:(\alpha^4,\alpha^2,1,\alpha^{-2},\alpha^{-4})$,
the $\Sym^2$ convolution bank, and the rank-4 Rankin bank
$c=\zeta\star\mathrm{sym2Bank}$.

\begin{theorem}[tensor-square decomposition]\label{thm:tensor}
The direct sum of the $\Sym^2$ clock system and the unit system carries the rank-4 Rankin
bank:
\[
  c_{\Sym^2\oplus\mathbf 1}\;=\;\zeta\star\mathrm{sym2Bank}\;=\;c .
\]
This is $\mathrm{std}\otimes\mathrm{std}=\Sym^2\oplus\mathbf 1$ at the bank level: the
closure law converts the direct-sum decomposition of the tensor square into the Euler-product
factorization $L(f\times f,s)=\zeta(s)\,L(\Sym^2\!f,s)$ of the $\Sym^2$ companion.
\end{theorem}
\begin{proof}
By Theorem~\ref{thm:bank}, $c_{\Sym^2\oplus\mathbf 1}=c_{\Sym^2}\star c_{\mathbf 1}$.  The
clock bank of $\Sym^2$ is the convolution bank (the coefficient identification of the
companion), and the unit bank is $\zeta$
(\texttt{bankArithmetic\_symClock\_eq\_sym2Bank}, \texttt{bankArithmetic\_unit\_eq\_zeta}).
(\texttt{tensorSquare\_bank}.)
\end{proof}

\begin{theorem}[plethysm: $\Sym^2\circ\Sym^2=\Sym^4\oplus\mathbf 1$]\label{thm:plethysm}
The $\Sym^2\!\circ\Sym^2$ weight system --- the six pairwise products
$w_iw_j$, $i\le j$, of the $\Sym^2$ clock --- decomposes as the $\Sym^4$ clock together with
one trivial channel, and consequently its bank is
\[
  c_{\Sym^2\circ\Sym^2}\;=\;\zeta\star c_{\Sym^4}.
\]
\end{theorem}
\begin{proof}
The six products are $\alpha^4,\alpha^2,1,1,\alpha^{-2},\alpha^{-4}$
(\texttt{sym2CompWeight\_products}); relabeling by the explicit bijection that sends the
second trivial channel to the unit summand exhibits the multiset as
$\Sym^4$-clock $\sqcup\{1\}$ (\texttt{sixSplit}, \texttt{sym2CompWeight\_decomposition}).
Apply relabeling invariance, Theorem~\ref{thm:bank}, the rank-uniform clock-bank
identification at $r=4$, and the unit bank.  (\texttt{plethysm\_bank}.)
\end{proof}

\begin{remark}[register and falsifiers]
Everything above is proven and machine-checked; nothing is conditional, and no automorphy is
consumed --- the two instances read only the compiled Satake seed of the eigenform.  The
classical content recovered: the Clebsch--Gordan/plethysm decompositions of $\mathrm{GL}(2)$
weight systems and the corresponding Euler-product factorizations, realized as instances of
one closure law rather than case-by-case identities.  A counterexample would be a pair of
weight systems whose union bank differs from the convolution of their banks at some prime
power, contradicting Theorem~\ref{thm:bank}; or a level at which the six $\Sym^2\circ\Sym^2$
products fail the $\Sym^4\oplus\mathbf 1$ split, contradicting
Theorem~\ref{thm:plethysm}.  Each is formally stated, so a counterexample would require a
kernel failure.
\end{remark}

\begin{thebibliography}{9}
\bibitem{GJ} S.~Gelbart and H.~Jacquet, \emph{A relation between automorphic representations
of $\mathrm{GL}(2)$ and $\mathrm{GL}(3)$}, Ann.\ Sci.\ \'ENS \textbf{11} (1978), 471--542.
\bibitem{KimShahidi} H.~Kim and F.~Shahidi, \emph{Functorial products for
$\mathrm{GL}(2)\times\mathrm{GL}(3)$ and the symmetric cube for $\mathrm{GL}(2)$}, Ann.\ of
Math.\ \textbf{155} (2002), 837--893.
\bibitem{Langlands} R.~P.~Langlands, \emph{Beyond endoscopy}, in: Contributions to automorphic
forms, geometry, and number theory, Johns Hopkins Univ.\ Press (2004), 611--697.
\end{thebibliography}

\end{document}
